\documentclass[10pt,a4paper]{amsart}
\usepackage[margin=19mm]{geometry}
\usepackage{amsmath,amssymb,mathtools}
\usepackage[scr=rsfs,cal=euler]{mathalfa}
\usepackage{xfrac}
\usepackage[unicode,hidelinks,bookmarks=false]{hyperref}
\newcommand{\ie}{i.e.\ }
\newcommand{\cf}{cf.\ }
\newcommand{\resp}{respectively\ }
\newcommand{\Subsection}{\subsection}
\providecommand{\nobreakdash}{\nobreak\hbox{-}}
\setlength{\emergencystretch}{2em}
\multlinegap0pt
\usepackage{bm}
\usepackage{bbm}
\usepackage{dsfont}
\usepackage{xspace}
\usepackage{cancel}
\usepackage{mathtools}
\usepackage{amsthm}
\makeatletter
\@ifundefined{thm}{\newtheorem{thm}{Theorem}}{}
\makeatother
\newcommand{\Om}{\Omega}
\newcommand{\ov}{\overline}
\title[Weighted Kernel Functions on Planar Domains]{Weighted Kernel Functions on Planar Domains}
\author{Aakanksha Jain, Kaushal Verma}
\date{Source: arXiv:2508.14016v1 (CC BY 4.0)}
\begin{document}
\maketitle

\noindent\emph{Source and licence.} Weighted Kernel Functions on Planar Domains. Version arXiv:2508.14016v1 (CC BY 4.0), distributed under the Creative Commons Attribution 4.0 International licence, as recorded in the arXiv licence metadata for this exact version and shown on the abstract page https://arxiv.org/abs/2508.14016v1. Full rights notice: licenses/arxiv-2508.14016-CC-BY-4.0.txt.\par\smallskip

\noindent\textit{Editorial scope.} Counted unit: the Thm whose source label is \texttt{None} in the article (source0.tex lines 439--457, source label \texttt{None}), delivered together with the article's own complete original proof of it (source0.tex lines 459--680, one \texttt{proof} environment of 222 lines). No mathematics is written, repaired or re-derived: statement and proof are the article's own text line for line. Required context retained uncounted: none. Every definition, display and label of those lines is retained unchanged. Omitted from this package and not redistributed: 1--438, 458--458, 681--1642. Nothing inside a retained excerpt is omitted or altered: every retained body is delivered exactly as the article's lines are written, so no omission receipt of any kind accompanies this package. Citations: the retained excerpts carry no citation command at all, so no bibliography entry of the article is transcribed and no reference is redistributed. Reference adaptation: none was needed and none was made; every reference inside the delivered unit resolves inside it, so no unresolved reference or citation is produced. Printed slips kept literal: no printed word, symbol, hypothesis or number of the counted statement or of its proof was altered. Layout and class: the article's own class is \texttt{amsart} with packages including inputenc, amsmath, amssymb, amsfonts, latexsym, amsthm, calligra, tcolorbox; the wrapper is the lane's pinned \texttt{amsart} editorial wrapper at 10pt on A4, which restates the theorem-like environments the retained text needs and adds the article's own macro definitions that the retained text needs (\texttt{\textbackslash Om}, \texttt{\textbackslash ov}), with extra wrapper packages (bm, bbm, dsfont, xspace, cancel, mathtools). The delivered numbering is the unit's own. Assets: no retained span contains a figure environment, a graphics-inclusion command, a PDF-page inclusion, a table or a TikZ picture; no image file is copied into this package, the package declares no asset dependency and no image file is redistributed with it. Licence: version arXiv:2508.14016v1 of this article is distributed under the Creative Commons Attribution 4.0 International licence, as recorded in the arXiv licence metadata for this exact version and shown on the abstract page https://arxiv.org/abs/2508.14016v1. A full rights notice is delivered at licenses/arxiv-2508.14016-CC-BY-4.0.txt. Characters: every retained excerpt is pure ASCII, so the delivered engine, XeLaTeX with its default Latin Modern text font, reports no missing character.
	\begin{thm}	
		There exists $k_0 \ge 1$ such that for all $k\geq k_0$ the following holds: 
		
		\medskip
		
		Let $w_j\in\Om$ be a sequence that converges to $a$. The functions $S_{\varphi_k}(\cdot,w_j)$ and $L_{\varphi_k}(\cdot,w_j)$ have distinct zeros in $\Omega$ for large enough $j$. The combined zeros $Z^i_k(w_j)$ of $S_{\varphi_k}(z,w_j)$ and $L_{\varphi_k}(\cdot, w_j)$ in $\overline{\Omega}$, become simple zeros, and it is possible to order them so that for each $i=1,\ldots, n-1$, there is a point $Z_k^i(a)\in\ov{\Om}$ such that $Z^i_k(w_j)$ tends to $Z^i_k(a)$ as $j\rightarrow\infty$. 
		
		\medskip
		
		The functions $S_{\varphi_k}(\cdot,a)$ and $L_{\varphi_k}(\cdot,a)$ have distinct zeros in $\Omega$. The combined number of zeros of $S_{\varphi_k}(\cdot, a)$ and $L_{\varphi_k}(\cdot,a)$ in $\Om$, plus the number of zeros of  $S_{\varphi_k}(\cdot,a)$ on $\partial\Om$ $($which are also the zeros of $L_{\varphi_k}(\cdot,a)$ on $\partial\Om)$, counting multiplicity, equals $n-1$, and are same as $Z_{k}^i(a)$. Furthermore, all $Z_{k}^i(a)$ are distinct for $i=1,\ldots, n-1$, and therefore the zeros are simple.
		
		\medskip
		
		Moreover, 
		\[
		\lim_{j,k\rightarrow\infty} Z^i_{k}(w_j) = Z^i(a)
		\]
		where $Z^i(a)$ are the zeros of $S(\cdot,a)$ (and $L(\cdot,a)$ as well).
	\end{thm}

	\begin{proof}
		Let $\gamma_i$ denote the $n$ boundary curves of $\Om$. For the sake of simplicity, we denote the boundary curves in such a manner that $a\in \gamma_n$. 
		Let $\Hat{\Om}$ denote the double of $\Om$, $R(z)$ denote the antiholomorphic involution on $\Hat{\Om}$ which fixes the boundary, and $\Tilde{\Om} = R(\Om)$ denote the reflection of $\Om$ in $\Hat{\Om}$ across the boundary. For each $i\neq n$, choose a smoothly bounded simply connected neighborhood $U_i$ of $Z^i(a)$ in $\Hat{\Om}$ and a point $b\in \Om$ such that 
		
		\begin{enumerate}
			\item[(i)] $\overline{U_i}$ are mutually disjoint
			\item[(ii)] $z\in U_i$ if and only if $R(z)\in U_i$
			\item[(iii)] $\ov{U_i}\cap ({\gamma_n \cup \{b\}}) = \emptyset$ 
			\item[(iv)] $S(z,b) \neq 0$ for all $z\in \overline{U_i}\cap \ov{\Om}$.
		\end{enumerate}
		
		\medskip
		
		Let $C_i$ denote the boundary curve of $U_i$. We claim that there exists a neighborhood $U$ of $a$ in $\overline{\Om}$ such that $S(z,w)\neq 0$ for all $z\in C_i\cap \overline{\Om}$ and $w\in \ov{U}$. If it were not true, then for every positive integer $m$, there would exist $w_m\in\ov{\Om}$ with $\vert w_m - a\vert \leq 1/m$ and $z_m\in C_i\cap \overline{\Om}$ for some $i$ such that $S(z_m,w_m) = 0$. By the pigeonhole principle, infinitely many $z_m$ lie on some $C_{i_0}$. So, there exists a subsequence $z_{m_r}$ with $z_{m_r}\rightarrow z_0\in C_{i_0}$ as $r\rightarrow\infty$. This implies that $S(z_0,a) = 0$ which is a contradiction.
		
		\medskip
		
		We know that
		\begin{equation}\label{thms}
			\lim_{k\rightarrow\infty} S_{\varphi_k}(z,w) = S(z,w)
			\quad \text{and}\quad
			\lim_{k\rightarrow\infty} L_{\varphi_k}(z,w) = L(z,w)
		\end{equation}
		locally uniformly on $(\ov{\Om}\times\ov{\Om})\setminus \{(z,z): z\in\partial\Om\}$ and $(\ov{\Om}\times\ov{\Om})\setminus \{(z,z): z\in \overline{\Om}\}\}$ respectively. Therefore, we can choose an integer $k_0 \ge 1$ such that for all $k\geq k_0$, 
		
		\begin{enumerate}
			\item[(i)] $S_{\varphi_k}(z,w)$ does not vanish for $z\in C_i\cap \overline{\Om}$, $w\in \ov{U}$
			\item[(ii)] $L_{\varphi_k}(R(z),w)$ does not vanish for $z\in C_i\cap (\Tilde{\Om}\cup \partial\Om)$, $w\in \ov{U}$
			\item[(iii)] $S_{\varphi_k}(z,b)$ does not vanish for $z\in \ov{U_i}\cap \overline{\Om}$
			\item[(iv)] $L_{\varphi_k}(R(z),b)$ does not vanish for $z\in \ov{U_i}\cap (\Tilde{\Om}\cup \partial\Om)$.
		\end{enumerate}
		
		\medskip
		
		For $k\geq k_0$ and $w\in \ov{U}$, define the functions $G_k(z,w)$ by
		\begin{equation}
		G_k(z,w)
		=
		\begin{cases}
			\mathcal{S}_{\varphi_k}(z,w)& z\in U_i\cap \overline{\Omega}
			\\
			\overline{\mathcal{L}_{\varphi_k}(R(z),w)}& z\in U_i\cap \Tilde{\Omega}
		\end{cases}
		\end{equation}
		where 
		\begin{equation}
		\mathcal{S}_{\varphi_k}(z,w) = \frac{S_{\varphi_k}(z,w)}{S_{\varphi_k}(z,b)}
		\quad
		\text{and}
		\quad
		\mathcal{L}_{\varphi_k}(z,w) = \frac{L_{\varphi_k}(z,w)}{L_{\varphi_k}(z,b)}.
		\end{equation}
		Recall that
		\[
			\overline{S_{\varphi_k}(z,w)} = \frac{1}{i\varphi_k(z)}L_{\varphi_k}(z,w)T(z),
			\quad\,\,
			w\in\Om, z\in\partial\Omega.
		\]
		This implies that
		\[
		\mathcal{S}_{\varphi_k}(z,w)
		=
		\frac{S_{\varphi_k}(z,w)}{S_{\varphi_k}(z,b)}
		=
		\overline{\left(\frac{L_{\varphi_k}(z,w)}{L_{\varphi_k}(z,b)}\right)}
		=
		\overline{\left(\frac{L_{\varphi_k}(R(z),w)}{L_{\varphi_k}(R(z),b)}\right)}
		=
		\overline{\mathcal{L}_{\varphi_k}(R(z),w)},
		\quad
		z\in\partial\Om.
		\]
		Therefore, the functions $G_k$ are well-defined and holomorphic on $U_i$ and extend smoothly to $C_i$. Similarly, we define the function $G(z,w)$ by
		\begin{equation}
		G(z,w)
		=
		\begin{cases}
			\mathcal{S}(z,w)& z\in U_i\cap \overline{\Omega}
			\\
			\overline{\mathcal{L}(R(z),w)}& z\in U_i\cap \Tilde{\Omega}
		\end{cases}
		\end{equation}
		where 
		\begin{equation}
		\mathcal{S}(z,w) = \frac{S(z,w)}{S(z,b)}
		\quad
		\text{and}
		\quad
		\mathcal{L}(z,w) = \frac{L(z,w)}{L(z,b)}.
		\end{equation}
		The function $G$ is well-defined, holomorphic on $U_i$ and extends smoothly to $C_i$. 
		Now, the function 
		\begin{equation}
		N_{ki}(w) = \int_{C_i} \frac{{G_k}'(z, w)}{G_k(z,w)} dz
		\end{equation}
		converges uniformly to 
		\begin{equation}
		N_i(w)=\int_{C_i}  \frac{{G}'(z, w)}{G(z,w)} dz
		\end{equation}
		for $w\in \overline{U}$ as $k\rightarrow\infty$. Here, $N_{ki}(w)$ and $N_i(w)$ are the number of zeros of $G_k(\cdot, w)$ and $G(\cdot,w)$ in $U_i$, respectively. We go back and choose a larger $k_0 \ge 1$, if needed, so that
		\[
		\vert N_{ki}(w) - N_i(w)\vert < 1/2
        \]
		for every $i,w\in \overline{U}$ and $k\geq k_0$.

        \medskip
		
		We shall now fix a $k\geq k_0$. Let $w_j\in\Omega$ be a sequence that converges to $a$. Assume, without loss of generality, that $w_j\in U$ for all $j$. The zeros of $G(\cdot, w_j)$ in $U_i$ are same as the zeros of $S(\cdot, w_j)$. So, there exists $j_0 \ge 1$ such that $N_i(w_j) = 1$ for all $i$ and $j\geq j_0$. Therefore, $N_{ki}(w_j) = 1$ for all $i$ and $j\geq j_0$. By the continuity of $N_{ki}(w)$ as a function of $w$, this implies that $N_{ki}(a) = 1$.
		
		\medskip
		
		The zeros of $G_k(\cdot,w_j)$ are the zeros of $S_{\varphi_k}(\cdot,w_j)$ or $L_{\varphi_k}(R(\cdot),w_j)$. The combined number of zeros of $S_{\varphi_k}(\cdot, w_j)$ and $L_{\varphi_k}(\cdot, w_j)$ in $\Om$, plus the number of zeros of $S_{\varphi_k}(\cdot,w_j)$ on $\partial\Om$ (which are also the zeros of $L_{\varphi_k}(\cdot,w_j)$ on $\partial\Om$), equals $n-1$. Since $N_{ki}(w_j) =1$ for all $i=1,\ldots,n-1$ and $j\geq j_0$, the functions $S_{\varphi_k}(\cdot,w_j)$ and $L_{\varphi_k}(\cdot,w_j)$ must have distinct zeros in $\Om$. Furthermore, the combined zeros $Z^i_k(w_j)$ of $S_{\varphi_k}(\cdot,w_j)$ and $L_{\varphi_k}(\cdot,w_j)$ in $\overline{\Om}$, are simple for all $j\geq j_0$. We order these zeros in such a manner that $Z^i_k(w_j)\in U_i$ for all $j\geq j_0$.
		
		\medskip
		
		Let $Z_i(G_k(\cdot,w_j))$ and $Z_i(G_k(\cdot,a))$ denote the zero of $G_k(\cdot,w_j)$ and $G_k(\cdot,a)$ in $U_i$ respectively. By an application of the residue theorem and (\ref{thms}),
		\[
		Z_i(G_k(\cdot,w_j))= \int_{C_i} \frac{z\,{G_k}'(z, w_j)}{G_k(z,w_j)} dz 
		\quad 
		\rightarrow
		\quad
		\int_{C_i}  \frac{z\,{G_k}'(z, a)}{G_k(z,a)} dz 
		=
		Z_i(G_k(\cdot,a))
        \]
		as $j\rightarrow\infty$.

    \medskip
        
		If $Z_i(G_k(\cdot,a))\in \ov{\Om}$, we let $Z^i_k(a) = Z_i(G_k(\cdot,a))$. Otherwise, we let $Z^i_k(a) = R(Z_i(G_k(\cdot,a)))$. Then, 
		\[
		\lvert Z^{i}_k(w_j)- Z^i_k(a) \vert 
		\leq 
		\vert Z_i(G_k(\cdot,w_j))- Z_i(G_k(\cdot,a)) \vert
		\rightarrow 0
		\quad
		\text{as }j\rightarrow\infty.
		\]
		Furthermore, $Z^i_k(a)\in U_i$ are distinct.
		
		\medskip
		
		Since the zeros of $G_k(\cdot,a)$ are the zeros of $S_{\varphi_k}(\cdot,a)$ or $L_{\varphi_k}(R(\cdot),a)$, the points $Z^i_k(a)$ are the zeros of $S_{\varphi_k}(\cdot,a)$ or $L_{\varphi_k}(\cdot,a)$ in $\overline{\Om}$. Let, if possible, $z_0\in\ov{\Om}$ be another zero of $S_{\varphi_k}(\cdot,a)$ or $L_{\varphi_k}(\cdot,a)$. Let us assume, without loss of generality, that $S_{\varphi_k}(z_0,a) = 0$, as a similar reasoning can be applied if $L_{\varphi_k}(z_0,a)=0$.  It is clear that $z_0\notin \overline{(U_i\cap \Om)}$ for any $i=1,\ldots,n-1$. If $z_0\in\Om$, choose a ball $B(z_0,\delta)\subset\subset\Om$ centered at $z_0$ of radius $\delta>0$ such that 
		
		\begin{itemize}
			\item[(i)] $\overline{B(z_0,\delta)}$ is disjoint from $\ov{U_i\cap \Om}$
			\item[(ii)] $S_{\varphi_k}(\cdot,a)$ is non-vanishing on $C=\partial B(z_0.\delta)$. 
		\end{itemize}
		
		\noindent Then $S_{\varphi_k}(\cdot,w_j)$ is also non-vanishing on $C$ for all large $j$. If it were not true, then for every $j \ge 1$, there would exist $z_j\in C$ such that $S_{\varphi_k}(z_j,w_j) = 0$. By passing to a convergent subsequence of $z_j$, we would obtain a zero of $S_{\varphi_k}(\cdot,a)$ on $C$, which is not possible. As $j\rightarrow\infty$,
		\[
		\int_{C}\frac{{S_{\varphi_k}}'(z,w_j)}{S_{\varphi_k}(z,w_j)} dz
		\quad
		\rightarrow
		\quad
		\int_{C}\frac{{S_{\varphi_k}}'(z,a)}{S_{\varphi_k}(z,a)} dz
		\quad
		\geq 1.
		\]
		So, there must exist a zero of $S_{\varphi_k}(\cdot,w_j)$ in $B(z_0,\delta)$ for all large $j$, hence a contradiction. If $z_0\in\partial\Om$, we will work with $G_k(\cdot,a)$ on an appropriate neighborhood of $z_0$, and arrive at a contradiction by a similar reasoning. Thus, we have proved that $Z^i_k(a)$ are the only zeros of $S_{\varphi_k}(\cdot,a)$ and $L_{\varphi_k}(\cdot,a)$. Since $N_{ki}(a) = 1$, it follows that $S_{\varphi_k}(\cdot,a)$ and $L_{\varphi_k}(\cdot,a)$ have distinct zeros in $\Om$. 
		
		\medskip
		
		We have also proved that the combined number of zeros of $S_{\varphi_k}(\cdot, a)$ and $L_{\varphi_k}(\cdot,a)$ in $\Om$, plus the number of zeros of  $S_{\varphi_k}(\cdot,a)$ on $\partial\Om$ $($which are also the zeros of $L_{\varphi_k}(\cdot,a)$ on $\partial\Om)$, counting multiplicity, equals $n-1$, and are same as $Z_{k}^i(a)$. Moreover, these zeros are simple.
		
		\medskip
		
		Finally, let $\epsilon>0$ be given. Using (\ref{thms}), we can choose $k_1\ge 1$ with $k_1\geq k_0$ such that
		\[
		\left\lvert
		\int_{C_i} z \; \frac{{G'_k}(z, w)}{G_k(z,w)} dz - \int_{C_i}  z \; \frac{G'(z, w)}{G(z,w)} dz 
		\right\rvert
		\,\,<\,\,
		\frac{\epsilon}{2}
		\quad
		\text{for all }i=1,\ldots, n-1,
		\,w\in\overline{U}
		\text{ and }
		k\geq k_1.
		\]
		Choose $j_1 \ge 1$ with $j_1\geq j_0$ such that
		\[
		\left\lvert
		\int_{C_i} z \; \frac{G'(z, w_j)}{G(z,w_j)} dz - \int_{C_i} z \; \frac{G'(z, a)}{G(z,a)} dz 
		\right\rvert 
		\,\,<\,\,
		\frac{\epsilon}{2}
		\quad
		\text{for all }i=1,\ldots, n-1
		\text{ and }
		j\geq j_1.
		\]
		Since $Z^i(a)\in\partial\Omega$, we have $\vert Z^i_k(w_j) - Z^i(a) \vert = \vert Z_i(G_k(\cdot,w_j)) - Z^i(a)\vert$. Therefore, for $j\geq j_1$ and $k\geq k_1$, we obtain using the residue theorem that
		\begin{multline*}
			\vert Z^i_k(w_j) - Z^i(a) \vert
			=
			\vert Z_i(G_k(\cdot,w_j)) - Z^i(a)\vert
			\\
			=
			\left\lvert
			\int_{C_i} z \; \frac{G'_k(z, w_j)}{G_k(z,w_j)} dz - \int_{C_i} z \; \frac{G'(z, a)}{G(z,a)} dz
			\right\rvert
			\\
			\leq
			\left\lvert
			\int_{C_i} z \; \frac{G'_k(z, w_j)}{G_k(z,w_j)} dz - \int_{C_i} z \;\frac{G'(z, w_j)}{G(z,w_j)} dz
			\right\rvert
			+
			\left\lvert
			\int_{C_i} z \;\frac{G'(z, w_j)}{G(z,w_j)} dz - \int_{C_i} z \; \frac{G'(z, a)}{G(z,a)} dz
			\right\rvert
			\\
			<
			\frac{\epsilon}{2} + \frac{\epsilon}{2} 
			= \epsilon.
		\end{multline*}
		Thus, 
		\begin{equation}
		\lim_{j,k\rightarrow\infty} Z^i_{k}(w_j) = Z^i(a)
		\end{equation}
		and this completes the proof of the theorem.		
	\end{proof}

\end{document}
